\chapter{Last-FM* Model Training and Evaluation}
\label{ch:lastfm-ewaluacja}

This chapter specifies the training objective, checkpoint selection and full-catalog validation used for the Last-FM* model of Chapter~\ref{ch:lastfm-proxy}.
Sampled validation guides iterative training; full-catalog validation is the primary ranking evaluation on the eligible catalog.

\section{Training Objective and Optimization}
\label{sec:lastfm-train}

For each training pair $i$ with label $y_i\in\{0,1\}$, the model produces a logit $s_i=s(u_i,X_i)$.
The sigmoid-transformed quantity
\begin{equation}
\label{eq:lastfm-phat}
\hat{p}_i=\sigma(s_i)
\end{equation}
is used inside the binary cross-entropy loss.
It is a training transform of the ranking score, not an externally calibrated probability of future interaction.

Training minimizes weighted binary cross-entropy,
\begin{equation}
\label{eq:lastfm-bce}
\mathcal{L}
=
-\frac{1}{N}\sum_i
\Bigl[
w_{+}\,y_i\log\hat{p}_i
+(1-y_i)\log(1-\hat{p}_i)
\Bigr],
\end{equation}
with positive-class weight
\begin{equation}
w_{+}=\frac{N_{\mathrm{neg}}}{N_{\mathrm{pos}}}.
\end{equation}
Under the frozen $4{:}1$ candidate construction of Section~\ref{sec:lastfm-negatives}, $w_{+}\approx 4$.
The positive-class weight balances the aggregate contribution of positive and negative training examples to the optimization objective.
This prevents the larger sampled-negative class from dominating the loss and encourages discrimination between observed interactions and sampled non-interactions.
Weighting does not make the classes intrinsically more separable.

Frozen optimization settings are:
Adam with learning rate $10^{-3}$ and weight decay $10^{-4}$;
ReduceLROnPlateau on sampled validation NDCG@20 (mode $\max$, factor $0.5$, patience $3$, threshold $10^{-4}$, minimum learning rate $6.25\times 10^{-5}$);
batch size $4096$;
minimum $15$ epochs and safety ceiling $120$ epochs (resolved training runs used an effective ceiling of $60$);
early-stopping patience $8$ with $\mathrm{min\_delta}=10^{-4}$ on sampled NDCG@20;
\texttt{float32} precision without automatic mixed precision;
no gradient clipping.

\section{Checkpoint Selection with Sampled NDCG@20}
\label{sec:lastfm-sampled-val}

Sampled validation provides a computationally efficient ranking signal during iterative training.
It is used for checkpoint selection, learning-rate scheduling and early stopping.
It is not the final ranking benchmark.

For each validation positive interaction of a user, the candidate list contains that positive item together with $20$ sampled negatives drawn from items outside the user's known history and known positives, subject to a frozen per-user cap of at most $10$ validation positives.
Relevance is binary on that list: the held-out positive is relevant; sampled negatives are non-relevant.
NDCG@20 is then computed on this sampled candidate universe.

\section{Full-Catalog Validation}
\label{sec:lastfm-fullcatalog}

Primary ranking evaluation uses the full eligible catalog
\begin{equation}
\label{eq:lastfm-Cu}
\mathcal{C}_u=I\setminus H_u^{(\mathrm{model\_train})},
\end{equation}
where $I$ is the catalog and $H_u^{(\mathrm{model\_train})}$ is the set of items observed for $u$ in \texttt{model\_train}.
Model-train seen items are masked; validation positives remain eligible; no sampled-negative approximation, popularity pruning or ANN pruning is applied.
Ties are broken deterministically by lower item identifier when scores coincide.

Let $P_u\subseteq\mathcal{C}_u$ denote the validation relevance set for user $u$.
For a ranked list of length $K$,
\begin{align}
\mathrm{DCG@}K
&=
\sum_{k=1}^{K}\frac{2^{r_k}-1}{\log_2(k+1)},
\label{eq:lastfm-dcg}\\
\mathrm{IDCG@}K
&=
\sum_{k=1}^{|P_u^{\ast}|}\frac{2^{1}-1}{\log_2(k+1)},
\label{eq:lastfm-idcg}\\
\mathrm{NDCG@}K
&=
\frac{\mathrm{DCG@}K}{\mathrm{IDCG@}K},
\label{eq:lastfm-ndcg}
\end{align}
where $r_k=1$ if the item at rank $k$ is relevant and $0$ otherwise, and $P_u^{\ast}$ is the ideal prefix of length $\min(K,|P_u|)$.
Recall@K is
\begin{equation}
\label{eq:lastfm-recall}
\mathrm{Recall@}K
=
\frac{|\{X\in\mathrm{Top}K(u)\}\cap P_u|}{|P_u|}.
\end{equation}
MRR is reported where applicable as the reciprocal rank of the first relevant item.

Sampled and full-catalog NDCG@20 apply the same ranking metric to different candidate universes.
The full-catalog protocol therefore represents a substantially more demanding ranking problem, and the numerical values are not directly comparable as measures of equal ranking difficulty.
